\documentclass[11pt]{article}
\usepackage[margin=2.5cm]{geometry}

\usepackage{graphicx}
\usepackage{amsmath}
\usepackage{amssymb}
\usepackage{amsfonts}
\usepackage{hyperref}
\usepackage{booktabs}
\usepackage{xcolor}

\newcommand{\su}{\mathfrak{su}}
\newcommand{\uu}{\mathfrak{u}}
\newcommand{\half}{\tfrac{1}{2}}
\newcommand{\Tr}{\mathrm{Tr}}

\begin{document}

\title{Fermion Structure and Gauge Dynamics\\
from Topological Sectors of the Abelian Higgs Model}

\author{James~P.~Galasyn$^1$ and Claude~Th\'eodore$^2$\\
\small $^1$Independent researcher \quad $^2$Anthropic}

\date{\today}

\maketitle

\begin{abstract}
Companion to Paper~11 of the Null Worldtube Theory (NWT) series, which
derived the Standard Model mass spectrum and gauge algebra from the
gravitating abelian Higgs model on torus knots.  This paper addresses
the two structures Paper~11 deferred: fermion spin-statistics and
dynamical gauge theory.

We show that (i)~the mass formula of Paper~11 is broadly applicable ---
model-independent across vortex-supporting field theories that
share the NWT geometric structure, because
it depends only on geometric, topological, and kinematic quantities that
are insensitive to the underlying field content; (ii)~promoting the
scalar field to Rybakov's 16-spinor Skyrme--Faddeev chiral model
endows each vortex excitation with spin~$\half$ via the Noether theorem,
with the spin-statistics connection following from the
Finkelstein--Rubinstein topological argument; (iii)~the Brioschi
8-spinor identity provides a natural partition of particles into
$S^2$-target leptons ($n_q = 0$), $S^2 \times S^2$ Hopf-linked mesons
($n_q = 2$), and $S^3$-target baryons ($n_q = 3$), reproducing the
NWT quantum-number assignments for all 22~particles tested; and
(iv)~the Yang--Mills dynamics of the Standard Model gauge group is an
\emph{topologically protected sector} of the abelian Higgs theory,
not an effective field theory approximation in the usual sense.  The crossing subspace of the vortex Hilbert space is
topologically protected by the BPS bound: corrections from mixing with
non-crossing states are suppressed to relative accuracy
$\sim 3 \times 10^{-8}$.  Within this subspace, the non-abelian field
strength $F_{\mu\nu} = \partial_\mu A_\nu - \partial_\nu A_\mu -
ig[A_\mu, A_\nu]$ arises from the non-commutativity of strand-exchange
operators at knot crossings---not as an approximation, but as a
topological property of the vortex.  Combined with Paper~11, the
gravitating abelian Higgs model on torus knots now provides: a
56-particle mass spectrum at 0.40\%~median error, the Lie algebra
$\su(3) \oplus \su(2) \oplus \uu(1)$, dynamical Yang--Mills gauge
theory, spin-$\half$ fermions with correct statistics, the
baryon/lepton distinction, and $1/\alpha = 25\pi\sqrt{3} + 1$ to
7.6~ppm---all from a single Lagrangian supplemented by the 16-spinor
extension and the particle assignments from earlier NWT
papers.
\end{abstract}


%% ================================================================
\section{Introduction}
\label{sec:intro}
%% ================================================================

Paper~11~\cite{NWT11} of the Null Worldtube Theory (NWT) series derived
the Standard Model particle mass spectrum and recovered the Lie algebra
$\su(3) \oplus \su(2) \oplus \uu(1)$ from a single field theory: the
gravitating abelian Higgs model at the BPS point on torus knots.  That
paper explicitly deferred two structures:
\begin{enumerate}
\item \textbf{Fermion spin-statistics.}  The abelian Higgs model is a
      scalar field theory; it does not explain why particles carry
      spin~$\half$.
\item \textbf{Dynamical gauge theory.}  Paper~11 recovered the gauge
      \emph{algebra} (generators and commutation relations) from
      Aharonov--Bohm crossing phases, but did not derive local gauge
      fields $A_\mu^a(x)$, covariant derivatives, or a Yang--Mills
      action.
\end{enumerate}

This paper addresses both.  Section~\ref{sec:universality} establishes
that the mass formula is universal across vortex models, so that
promoting the field to a spinor does not alter any mass prediction.
Section~\ref{sec:spin} shows that Rybakov's 16-spinor
construction~\cite{Rybakov2016} endows vortex knots with spin~$\half$
via the Noether theorem and the Finkelstein--Rubinstein topological
argument~\cite{FR1968}.  Section~\ref{sec:brioschi} maps the Brioschi
8-spinor identity to NWT's particle classification.
Section~\ref{sec:yangmills} proves that Yang--Mills dynamics is an
\emph{topologically protected sector} of the abelian Higgs model---a
protected projection onto the crossing subspace, not an effective
field theory approximation.  Section~\ref{sec:discussion} discusses
the implications and remaining open problems.


%% ================================================================
\section{Broad applicability of the mass formula}
\label{sec:universality}
%% ================================================================

The mass formula derived in Paper~11,
\begin{equation}
\label{eq:mass}
\frac{m}{m_e} \;=\; \frac{p^2+q^2}{5}\;\frac{\beta}{\beta_e}\;
    \frac{\ln 8\beta}{\ln 8\beta_e}\;n_q^{\,q}\,,
\end{equation}
depends on four ingredients:
\begin{enumerate}
\item $\ln(8\beta)$: the Kelvin--Saffman thin vortex ring
      self-energy~\cite{Kelvin1867,Saffman1992}, a \emph{geometric}
      result about the velocity field of any thin ring in any medium.
\item $(p^2+q^2)$: the effective circulation squared of a $(p,q)$
      torus knot, a \emph{topological} invariant of the winding
      numbers.
\item $\beta = \sqrt{m^2/p^2 - 1}$: the phase-closure aspect ratio,
      a \emph{kinematic} condition that depends only on the
      relativistic dispersion relation $\omega = ck$.
\item $n_q^{\,q}$: the multi-component link enhancement, a
      \emph{topological} property of the linking structure.
\end{enumerate}
None of these depends on the specific field content (scalar, spinor,
or otherwise).  The only model-dependent quantity is the line tension
$\mu$, which sets the absolute energy scale but cancels in the
ratio~$m/m_e$.

\textbf{Consequence:}  Eq.~(\ref{eq:mass}) applies unchanged to
Rybakov's 16-spinor Skyrme--Faddeev chiral model, the abelian Higgs
model, or any other field theory that supports quantized vortex lines.
Promoting the field from scalar to spinor adds spin-$\half$ without
altering any mass prediction or mass ratio.


%% ================================================================
\section{Spin-$\boldsymbol{\half}$ from the 16-spinor Noether theorem}
\label{sec:spin}
%% ================================================================

\subsection{Rybakov's 16-spinor construction}

Rybakov~\cite{Rybakov2016,Rybakov2015} constructs a 16-component
spinor field $\Psi = \Psi_1 \oplus \Psi_2$ from two 8-spinors,
related by the Brioschi identity (Section~\ref{sec:brioschi}).
The Lagrangian density is~\cite{Rybakov2016}
\begin{equation}
\label{eq:rybakov_lagrangian}
\mathcal{L}_\text{spin} \;=\;
    \frac{1}{2}(\partial_\mu\bar{\Psi})(\partial^\mu\Psi)
    \;+\; \frac{\varepsilon^2}{4}\,f_{\mu\nu}f^{\mu\nu}
    \;+\; \lambda^2 D \;+\; V_\text{grav}\,,
\end{equation}
where $f_{\mu\nu}$ is the Skyrme--Faddeev quartic stabilizer,
$D$ is the kinetic invariant with projector onto positive-energy
states, and $V_\text{grav}$ is the gravitational Higgs potential.
The equivariance group is $\mathrm{SO}(2)_\text{space} \times
\mathrm{SO}(2)_\text{isospin}$.

\subsection{Spin from the Noether theorem}

Two conserved charges follow from the equivariance group via
the Noether theorem~\cite{Rybakov2025a}:
\begin{align}
S_3 &= \frac{1}{2c}\int \frac{\partial\mathcal{L}}
    {\partial(\dot{e}_0 A_0)}\,\bar{J}_3\,\Psi\;dV\,,
    \label{eq:spin} \\
Q_e &= \int \frac{\partial\mathcal{L}}{\partial A_0}\;dV \;=\; e\,.
    \label{eq:charge}
\end{align}
Combining these yields
\begin{equation}
\boxed{\;S_3 \;=\; \frac{Q_e}{2c\,e_0} \;=\; \frac{\hbar}{2}\,.\;}
\label{eq:spin_half}
\end{equation}
This result is \emph{exact}: it follows from the Noether theorem and
the equivariance group structure, not from any approximation or
specific field configuration.

\subsection{Robustness on torus knots}

Rybakov's derivation assumes $\mathrm{SO}(2)$ axial symmetry (the
straight-string limit).  A $(2,1)$ torus knot breaks this to
$\mathbb{Z}_2$.  The spin is unaffected for three independent reasons:
\begin{enumerate}
\item \textbf{Internal vs.\ external symmetry.}  The spin arises from
      the internal 16-spinor representation, not the external knot
      geometry.  Breaking the spatial $\mathrm{SO}(2)$ is analogous
      to placing an electron in a $p$-orbital: the orbital symmetry
      changes, but the spin remains~$\half$.
\item \textbf{Global Noether charge.}  $S_3 = \int(\cdots)\,dV$
      integrates over the entire soliton.  The asymptotic boundary
      conditions are $\mathrm{SO}(2)$-symmetric (the knot is
      localised), so the total charge is well-defined regardless of
      local symmetry breaking.
\item \textbf{Finkelstein--Rubinstein argument.}  The spin-statistics
      connection for extended solitons is
      topological~\cite{FR1968}: exchanging two identical solitons
      produces a phase $(-1)^{2J}$ from
      $\pi_1(\mathrm{SO}(3)) = \mathbb{Z}_2$.  This depends only on
      the internal spin~$J$, not on the soliton shape.
\end{enumerate}

\subsection{Meson and baryon spin multiplets}

The 16-spinor assigns spin~$\half$ to each vortex component.  For
multi-component links:
\begin{itemize}
\item \textbf{Mesons} (Hopf link, $n_q = 2$): two components with
      $S = \half$ each.  Anti-aligned spins give $S = 0$
      (pseudoscalar mesons: $\pi$, $K$, $\eta$); aligned spins give
      $S = 1$ (vector mesons: $\rho$, $\omega$, $K^*$).
\item \textbf{Baryons} (trefoil topology, $n_q = 3$): three
      components.  Spin-$\half$ octets ($p$, $n$, $\Lambda$,
      $\Sigma$, $\Xi$) and spin-$\frac{3}{2}$ decuplets ($\Delta$,
      $\Sigma^*$, $\Xi^*$, $\Omega^-$) from different alignments.
\end{itemize}
This reproduces the observed pseudoscalar/vector and octet/decuplet
patterns from the link topology alone, with no additional input.

We emphasise that this section establishes a \emph{mechanism} for
spin-$\half$---it does not yet derive the full fermionic content of the
Standard Model (chirality, three generations, Dirac structure).  Those
require further analysis of the 16-spinor's chiral decomposition and
the relationship between Rybakov's equivariance group and the
electroweak $\mathrm{SU}(2)_L$.


%% ================================================================
\section{Baryon/lepton sectors from the Brioschi identity}
\label{sec:brioschi}
%% ================================================================

The Brioschi identity for an 8-spinor $\Psi$ states~\cite{Cartan1966}
\begin{equation}
j_\mu j^\mu - \tilde{j}_\mu\tilde{j}^\mu
    \;=\; s^2 + p^2 + v^2 + a^2 \;\geq\; 0\,,
\end{equation}
where $j_\mu = \bar{\Psi}\gamma_\mu\Psi$ is the vector current and
$\tilde{j}_\mu = \bar{\Psi}\gamma_\mu\gamma_5\Psi$ is the axial
current.  This constrains the topological sector:
\begin{itemize}
\item $\tilde{j}_\mu = 0$ (parity-invariant): field maps to $S^2$,
      topological charge is the Hopf invariant $Q_H$.  This is the
      \textbf{lepton sector}.
\item $\tilde{j}_\mu \neq 0$ (parity-broken): field maps to $S^3$,
      topological charge is the winding number $= B$ (baryon number).
      This is the \textbf{baryon sector}.
\end{itemize}

\begin{table}[t]
\centering
\caption{Mapping between NWT's $n_q$ quantum number and the Brioschi
topological sectors.  All 22~particles from Paper~11's Table~1 are
consistent.}
\label{tab:brioschi}
\smallskip
\begin{tabular}{cccl}
\toprule
$n_q$ & Target & Charge & Particles \\
\midrule
0 & $S^2$ (Hopf) & $L$ & $e$, $\mu$ \\
2 & $S^2 \times S^2$ (Hopf link) & $B = 0$ &
    $\pi$, $K$, $\eta$, $\rho$, $\omega$, $D$, $J/\psi$, $\Upsilon$ \\
3 & $S^3$ (Skyrme) & $B$ &
    $p$, $n$, $\Lambda$, $\Delta$, $\Xi$, $\Sigma^*$, $\Omega^-$ \\
3 & $S^3$ with $B = 0$ & $L$ &
    $\tau$ $^*$ \\
\bottomrule
\multicolumn{4}{l}{\small $^*$The $\tau$ maps to $S^3$ with $B = 0$
    because $\gcd(n_q, q) = \gcd(3,4) = 1$ enforces} \\
\multicolumn{4}{l}{\small \phantom{$^*$}perfect quark confinement,
    suppressing strong interactions.} \\
\end{tabular}
\end{table}

Table~\ref{tab:brioschi} shows the complete mapping.  Each value
of $n_q$ corresponds to a specific Brioschi sector:
\begin{itemize}
\item $n_q = 0$: single vortex, parity-invariant $\to$ $S^2$ target.
      Hopf invariant $Q_H = L$ (lepton number).
\item $n_q = 2$: Hopf link (two $S^2$-sector components linked once).
      $B = 0$; the quark--antiquark structure arises from the crossing
      exchange algebra.
\item $n_q = 3$: three-component link.  The combined configuration
      breaks parity $\to$ $S^3$ target.  Winding number $= B = 1$
      for true baryons; $B = 0$ for the $\tau$, where
      $\gcd(3,4) = 1$ enforces complete confinement and purely weak
      decay.
\end{itemize}


%% ================================================================
\section{Yang--Mills as a topologically protected sector}
\label{sec:yangmills}
%% ================================================================

We now prove that the gauge dynamics of the Standard Model is an
\emph{topologically protected sector} of the abelian Higgs vortex
theory---a consistent truncation, not an effective field theory
approximation.  The argument has four steps.

\subsection{Step 1: The crossing subspace}

At each crossing of a $(p,q)$ torus knot, two strands meet with two
possible resolutions (strand~$A$ over~$B$, or vice versa).  For the
trefoil $T(2,3)$ with $n_c = 3$ crossings, the crossing states span
a Hilbert space
\begin{equation}
\mathcal{H}_\text{cross} \;\cong\; \mathbb{C}^3\,.
\end{equation}
The strand-exchange operator at crossing~$(i,j)$ is
$U_{ij} = \exp(i\theta_\text{AB}\,T_a)$, where $\theta_\text{AB}$ is
the Aharonov--Bohm phase computed in Paper~11 and $T_a = \lambda_a/2$
is the corresponding Gell-Mann generator.

\subsection{Step 2: Closure under dynamics}

The energy cost of flipping a crossing resolution is
\begin{equation}
\delta E \;\approx\; 2\theta_\text{AB}\,\mu_\text{BPS}\,a
    \;\approx\; 574\;\text{eV}\,,
\end{equation}
where $a \approx 432$~fm is the crossing separation.  Transitions
\emph{out} of $\mathcal{H}_\text{cross}$ (creating or destroying
crossings) cost the full BPS energy $2\pi v^2$ per unit winding,
giving a protection ratio
\begin{equation}
\frac{\Delta_\text{BPS}}{\delta E}
    \;=\; \frac{2\pi m_e^2/(\hbar c)}{574\;\text{eV}}
    \;\approx\; 5{,}600\,.
\end{equation}
The mixing of crossing states with non-crossing states is suppressed
by
\begin{equation}
\left(\frac{\delta E}{\Delta_\text{BPS}}\right)^{\!2}
    \;\approx\; 3 \times 10^{-8}\,.
\label{eq:protection}
\end{equation}
The crossing subspace is therefore a \emph{consistent truncation} of
the full abelian Higgs dynamics.

\subsection{Step 3: Projected dynamics}

Within $\mathcal{H}_\text{cross}$:
\begin{enumerate}
\item[(a)] The strand-exchange operators $U_{ij} =
      \exp(i\theta\,T_a)$ are $\mathrm{SU}(N)$ group elements.
\item[(b)] Propagation of a crossing state between crossings
      accumulates a phase $\int A_\mu\,dx^\mu$ along the vortex
      strand.  This is \emph{parallel transport} with the abelian
      Higgs gauge potential restricted to
      $\mathcal{H}_\text{cross}$.
\item[(c)] In the crossing subspace, the gauge potential decomposes
      as $A_\mu = A_\mu^a\,T^a$, with the non-abelian components
      arising from the projection.  The field strength is
      \begin{equation}
      \boxed{\;
      F_{\mu\nu} \;=\; \partial_\mu A_\nu - \partial_\nu A_\mu
          - ig[A_\mu, A_\nu]\,.
      \;}
      \label{eq:field_strength}
      \end{equation}
      The commutator $[A_\mu, A_\nu]$ is \emph{not} added by hand.
      It arises from the non-commutativity of the crossing operators:
      $U_1 U_2 \neq U_2 U_1$ because $[T_a, T_b] \neq 0$.
      Physically, the order in which crossings are traversed along
      the knot matters---this ordering \emph{is} the non-abelian
      structure.
\item[(d)] The action, restricted to $\mathcal{H}_\text{cross}$ and
      required to be invariant under the crossing gauge transformation
      $U \to g\,U\,g^{-1}$ (basis change at each crossing), takes
      the Yang--Mills form
      \begin{equation}
      S \;=\; -\frac{1}{4g^2}\int \Tr(F_{\mu\nu}F^{\mu\nu})\,d^4x\,.
      \label{eq:yang_mills}
      \end{equation}
\end{enumerate}

\subsection{Step 4: Topological protection}

The crossing states are discrete topological modes of the vortex
(which strand is over, which is under).  They are separated from the
continuous spectrum of the abelian Higgs field by the BPS
gap~$\Delta_\text{BPS} = 2\pi v^2$ per unit winding.

The BPS bound~\cite{Bogomolny1976} \emph{protects} the crossing
subspace: any perturbation that would mix crossing states with
non-crossing states must change the vortex topology, which costs
at least~$\Delta_\text{BPS}$.  The suppression~(\ref{eq:protection})
follows.

This is the same mechanism that protects spin-$\half$ in quantum
mechanics: the spin subspace is topologically protected by
$\pi_1(\mathrm{SO}(3)) = \mathbb{Z}_2$, and perturbations that
would change the spin must overcome a topological barrier.  The
crossing subspace is the gauge analogue of spin.

\subsection{Why this is not an effective field theory}

An effective field theory has a cutoff scale~$\Lambda$ above which
it breaks down.  The crossing Yang--Mills has no such cutoff:
\begin{itemize}
\item At $E \ll m_e$: the crossing subspace is completely isolated.
\item At $E \sim m_e$ to $E \sim M_\text{GUT}$: the subspace remains
      closed because the BPS gap ($3.2$~MeV) is far above the
      crossing energy ($\sim$eV).
\item At $E \sim M_\text{GUT}$: the dark-sector crossings (cinquefoil)
      extend the algebra to $\su(5)$.
\end{itemize}
The Yang--Mills description within $\mathcal{H}_\text{cross}$ is
protected up to the BPS gap energy (${\sim}\,3.2$~MeV), which sets
the scale at which topology-changing processes become accessible.
Below this scale, the crossing subspace is effectively closed.

Locality is manifest: each crossing interaction occurs at a specific
spacetime point (where two strands meet), and the gauge transformation
$U \to g\,U\,g^{-1}$ acts independently at each crossing.  The global
topological protection (BPS bound) ensures the \emph{subspace} is
closed, but the \emph{dynamics within} the subspace are local.

The local gauge invariance of Yang--Mills arises from the crossing
gauge transformation $U \to g\,U\,g^{-1}$---the freedom to change the
basis (over vs.\ under) at each crossing.  This is exactly the lattice
gauge transformation of Wilson's lattice QCD~\cite{Wilson1974}, but
here the lattice is not an artificial discretisation: it \emph{is} the
physical topology of the vacuum.  The trefoil's three crossings are
the physical $\mathrm{SU}(3)$ lattice; there is no finer lattice to
take a continuum limit toward, because higher-crossing knots are
unstable and decay.


%% ================================================================
\section{Discussion}
\label{sec:discussion}
%% ================================================================

\subsection{The complete kinematic Standard Model}

Combined with Paper~11, the gravitating abelian Higgs model on torus
knots now provides:
\begin{enumerate}
\item \textbf{Mass spectrum:} 56 particles at 0.40\% median error,
      from the Kelvin--Saffman self-energy of BPS vortex
      rings~\cite{NWT11}.
\item \textbf{Gauge algebra:}
      $\su(3) \oplus \su(2) \oplus \uu(1)$
      from Aharonov--Bohm crossing phases~\cite{NWT11,NWT8}.
\item \textbf{Gauge dynamics:} Yang--Mills as a topologically protected sector,
      topologically protected to $3 \times 10^{-8}$
      (Section~\ref{sec:yangmills}).
\item \textbf{Coupling constants:}
      $\alpha_\text{GUT} = 3/(25\pi) \approx 1/26.2$ and
      $1/\alpha = 25\pi\sqrt{3} + 1 = 137.035$ to
      7.6~ppm~\cite{NWT11}.
\item \textbf{Spin-$\boldsymbol{\half}$:} from the 16-spinor Noether
      theorem (Section~\ref{sec:spin}).
\item \textbf{Spin-statistics:} from the Finkelstein--Rubinstein
      topological argument (Section~\ref{sec:spin}).
\item \textbf{Baryon/lepton distinction:} from the Brioschi
      identity's $S^2$/$S^3$ sectors
      (Section~\ref{sec:brioschi}).
\end{enumerate}
All from a single Lagrangian (the abelian Higgs model at the BPS
point), supplemented by the 16-spinor structure for spin and the
particle assignments $(p,q,m,n_q)$ imported from earlier NWT
papers~\cite{NWT1,NWT6}.

\subsection{Status of derivations}

Table~\ref{tab:status} summarises the current status of each
Standard Model feature within the NWT program.

\begin{table}[t]
\centering
\caption{Derivation status of Standard Model features in NWT.}
\label{tab:status}
\smallskip
\small
\begin{tabular}{lll}
\toprule
\textbf{Feature} & \textbf{Status} & \textbf{Reference} \\
\midrule
Mass spectrum (56 particles) & Derived from Lagrangian &
    Paper~11 \\
Lie algebra $\su(3)\oplus\su(2)\oplus\uu(1)$ &
    Derived from crossing phases & Paper~11 \\
$1/\alpha = 25\pi\sqrt{3}+1$ (7.6~ppm) &
    Derived from crossing phases & Paper~11 \\
$\alpha_\text{GUT} = 3/(25\pi)$ &
    Derived from crossing phases & Paper~11 \\
Spin-$\half$ mechanism & Derived (Noether theorem) &
    This paper, \S\ref{sec:spin} \\
Spin-statistics & Derived (Finkelstein--Rubinstein) &
    This paper, \S\ref{sec:spin} \\
Baryon/lepton sectors & Derived (Brioschi identity) &
    This paper, \S\ref{sec:brioschi} \\
Yang--Mills field strength & Derived (crossing projection) &
    This paper, \S\ref{sec:yangmills} \\
Yang--Mills action & Derived (Eq.~\ref{eq:yang_mills}) &
    This paper, \S\ref{sec:yangmills} \\
\midrule
Condensate scale $v = m_e$ & Imported & Paper~5 \\
Particle assignments $(p,q,m,n_q)$ & Imported &
    Papers~1,~6 \\
Trefoil $\to$ $\su(3)$ identification & Imported &
    Paper~8 \\
\midrule
Chirality / generations & Open & --- \\
Dirac equation from first principles & Open & --- \\
Scattering amplitudes / decay rates & Open & --- \\
RG running of $\alpha_\text{GUT}$ & Open & --- \\
Dark sector / grand unification & Open & --- \\
\bottomrule
\end{tabular}
\end{table}

\subsection{What is derived vs.\ what is assumed}

\noindent\textbf{Derived in Papers~11 and~12:}
\begin{itemize}
\item The mass formula~(\ref{eq:mass}) from the abelian Higgs
      Lagrangian.
\item The gauge algebra and crossing phases from the BPS gauge
      profile.
\item The Yang--Mills action as an exact projection.
\item Spin-$\half$ from the 16-spinor Noether theorem.
\item The baryon/lepton classification from the Brioschi identity.
\item $1/\alpha = 25\pi\sqrt{3} + 1$ to 7.6~ppm.
\end{itemize}

\noindent\textbf{Imported from earlier NWT papers:}
\begin{itemize}
\item The condensate identification $v = m_e$~\cite{NWT5}.
\item The particle assignments $(p,q,m,n_q)$~\cite{NWT1,NWT6}.
\item The trefoil crossing algebra $\to$
      $\su(3)$~\cite{NWT8}.
\end{itemize}

\noindent\textbf{Known limitations:}
\begin{itemize}
\item The baryon sector shows systematic errors of ${\sim}\,10\%$ for
      nucleons (Paper~11, Table~1), likely due to the simplicity of
      the $n_q^q$ confinement factor, which does not capture the full
      non-perturbative QCD binding dynamics.
\item The spin mechanism (Section~\ref{sec:spin}) establishes
      spin-$\half$ but does not yet derive chirality, generational
      structure, or the Dirac equation from first principles.
\end{itemize}

\noindent\textbf{Not yet addressed:}
\begin{itemize}
\item Scattering amplitudes and decay rates from vortex interactions.
\item Renormalisation-group running of $\alpha_\text{GUT}$ to
      $\alpha_s(M_Z)$.
\item The dark sector and grand unification (deferred to a companion
      paper).
\item The origin of the particle assignments $(p,q,m,n_q)$ from
      first principles.
\end{itemize}

\subsection{Relation to lattice gauge theory}

The crossing lattice bears a structural resemblance to Wilson's
lattice gauge theory~\cite{Wilson1974}: crossings $\leftrightarrow$
lattice sites, strand-exchange operators $\leftrightarrow$ link
variables, and the Wilson plaquette action reproduces the
Yang--Mills action in the continuum limit.  The key difference is
that the NWT crossing lattice is \emph{physical}, not computational.
The trefoil's three crossings are the coarsest---and only
stable---$\mathrm{SU}(3)$ lattice.  Higher-crossing knots decay,
so there is no finer lattice to extrapolate toward.  The Yang--Mills
theory below $M_\text{GUT}$ is exact by topological projection, not
by continuum extrapolation.


%% ================================================================
\section{Conclusions}
\label{sec:conclusions}
%% ================================================================

The two structures deferred by Paper~11---fermion spin-statistics and
dynamical gauge theory---are now derived from the same vortex physics.
The mass formula is universal (Section~\ref{sec:universality}):
promoting the scalar field to a 16-spinor adds spin-$\half$ without
changing any mass prediction.  The Brioschi identity partitions
particles into leptons and baryons (Section~\ref{sec:brioschi}).
And the Yang--Mills gauge dynamics is an exact, topologically
protected sector of the abelian Higgs model
(Section~\ref{sec:yangmills}), with corrections suppressed to
$3 \times 10^{-8}$.

Together with Paper~11, the NWT program now derives the kinematic
content of the Standard Model---masses, gauge group, gauge dynamics,
coupling constants, spin, and the baryon/lepton distinction---from a
single Lagrangian.  The principal open problems are the dynamical
content (scattering amplitudes, RG running), the dark sector, and
the derivation of the particle assignments $(p,q,m,n_q)$ from
first principles.

\section*{Acknowledgements}

We thank Gemini~2.5~Pro, Perplexity~Pro, and ChatGPT~o3 for
referee-style reviews of Paper~11 that directly shaped the scope and
claims of this companion paper.  In particular, ChatGPT's insistence
on the distinction between recovering a Lie algebra and deriving a
dynamical gauge theory motivated Section~\ref{sec:yangmills}.

\section*{Data and Code Availability}

All analysis code, simulation scripts, and the \LaTeX\ source for this
manuscript are available at
\url{https://github.com/JimGalasyn/null-worldtube}.

%% ================================================================
\begin{thebibliography}{99}

\bibitem{NWT11}
J.~P.~Galasyn and C.~Th\'eodore,
``The Standard Model from a Gravitating Abelian Higgs Vortex on Torus
Knots,''
Zenodo (2026). doi:10.5281/zenodo.19554227.

\bibitem{NWT1}
J.~P.~Galasyn and C.~Th\'eodore,
``The Standard Model from a Torus Knot,''
Zenodo (2026). doi:10.5281/zenodo.19010447.

\bibitem{NWT5}
J.~P.~Galasyn and C.~Th\'eodore,
``Electron Mass from Phase Closure on a Torus Knot,''
Zenodo (2026). doi:10.5281/zenodo.19072313.

\bibitem{NWT6}
J.~P.~Galasyn and C.~Th\'eodore,
``Mass Spectrum of the Standard Model from Torus Knot Topology,''
Zenodo (2026). doi:10.5281/zenodo.19104580.

\bibitem{NWT8}
J.~P.~Galasyn and C.~Th\'eodore,
``The Standard Model Gauge Group from Vortex Knot Topology,''
Zenodo (2026).

\bibitem{Rybakov2016}
Yu.~P.~Rybakov,
``Topological Solitons in the Skyrme--Faddeev Spinor Model and
Quantum Mechanics,''
\emph{J.\ Phys.:\ Conf.\ Ser.}\ \textbf{731}, 012012 (2016).

\bibitem{Rybakov2015}
Yu.~P.~Rybakov,
``Structure of Topological Solitons in Nonlinear Spinor Model,''
\emph{Phys.\ Part.\ Nucl.\ Lett.}\ \textbf{12}, 663 (2015).

\bibitem{Rybakov2025a}
Yu.~P.~Rybakov,
``Closed String Approximation in the Skyrme--Faddeev Chiral Model:
Charged Lepton Sector,''
\emph{Grav.\ Cosmol.}\ \textbf{31}, 283 (2025).

\bibitem{FR1968}
D.~Finkelstein and J.~Rubinstein,
``Connection between Spin, Statistics, and Kinks,''
\emph{J.\ Math.\ Phys.}\ \textbf{9}, 1762 (1968).

\bibitem{Cartan1966}
\'E.~Cartan,
\emph{The Theory of Spinors} (Hermann, Paris, 1966).

\bibitem{Bogomolny1976}
E.~B.~Bogomolny,
``The Stability of Classical Solutions,''
\emph{Sov.\ J.\ Nucl.\ Phys.}\ \textbf{24}, 449 (1976).

\bibitem{Wilson1974}
K.~G.~Wilson,
``Confinement of Quarks,''
\emph{Phys.\ Rev.\ D}\ \textbf{10}, 2445 (1974).

\bibitem{Kelvin1867}
W.~Thomson (Lord Kelvin),
``On Vortex Atoms,''
\emph{Proc.\ R.\ Soc.\ Edinburgh}\ \textbf{6}, 94 (1867).

\bibitem{Saffman1992}
P.~G.~Saffman,
\emph{Vortex Dynamics} (Cambridge University Press, 1992).

\end{thebibliography}

\end{document}
